\chapter{附录A 习题参考答案}
\label{app:answers}

本附录按章节和题号给出客观题答案、简答题评分要点与实践题验收清单。教师宜在布置作业并完成课堂讨论后开放本附录；评分要点允许学生采用不同实现，只要证据能够复核。

\section{第1章}
\begin{description}
\item[1--4（评分要点）] 说明平台是数据、模型、业务流程和责任链的组合；区分标准与风险；分别给出监测上行和控制下行；WaterML 2.0 是数据交换模型而非通信协议。每题答出概念边界、一个水利例子和一个工程后果各得分。
\item[5（计算）] $28\times200\times60\times24=8{,}064{,}000$ 字节，约 8.06 MB/日（十进制）。还应估计协议头、时间戳、重传、索引、备份、副本和峰值余量。
\item[6--8（评分要点）] 审计证据至少含原始观测、质量检查、模型/规则版本、预警事件、确认人、工单回执和时间戳；降级方案需定义触发、保留能力、恢复和告警；模拟建议与真实控制之间要有权限、审批、设备状态和人工确认边界。
\end{description}

\section{第2章}
\begin{description}
\item[1--3（评分要点）] 生命周期包含退役以处理数据、权限和资产；稳定接口项目可选阶段门禁明确的增量/瀑布变体并说明变更风险；原型用于验证假设，生产还需安全、性能、可运维和数据迁移证据。
\item[4] 干系方可含管理者、调度员、用水户、运维员、审批人和审计员；每类写出目标、输入和冲突。
\item[5--6] MoSCoW 应给出五条需求的 Must/Should/Could/Won't、依赖与理由；“快速”改为可测的 P95 响应、数据范围、并发量和验收时间窗。
\item[7--9] 异常路径从过期计划回退到编辑或重新申报；DFD 上下层数据流名称和数量平衡；数据流只能表示数据对象，字段需有类型、单位、取值域、来源、校验和责任人。
\item[10--12] 决策表列出条件组合和动作且无重叠；SRS 至少含范围、角色、功能/非功能需求、接口、数据、验收和变更规则；追踪矩阵把需求连到用例、设计、测试和证据。
\end{description}

\section{第3章}
\begin{description}
\item[1--4（评分要点）] 模块化以低耦合高内聚降低复杂度；MVC、微服务、事件驱动按部署边界和一致性需求选用；抽象稳定变化点，信息隐藏隔离实现；架构评估应以性能、安全、可用性和演进风险选择方法。
\item[5--8] 大数据平台采用分层/事件驱动并说明吞吐、成本和运维；分层、C/S、分布式分别说明边界和适用系统；权衡记录功能、非功能、成本、风险和验证手段；并发优化需负载均衡、缓存失效、限流和观测指标。
\item[9--12] 架构设计应画层次、模块、接口、数据和部署；UML 用例/类/时序需同一场景一致；清源水库分析要写决策、权衡和改进；微服务按业务边界拆分，事件用 Kafka，数据按服务负责且以审计/幂等约束协同。
\end{description}

\section{第4章}
\subsection*{客观题答案}
1 A；2 A；3 错；4 B；5 A；6 对。
\begin{description}
\item[1--5（评分要点）] DOM/CSSOM 构成渲染树并由脚本触发样式与布局；三类卡片使用语义元素、可访问名称和 Grid 断点；Promise 与 async/await 都处理错误和取消；路由表写路径、参数、权限，Pinia 只放跨页面状态；性能诊断提交瀑布、Long Task、布局次数和修复前后证据。
\end{description}

\section{第5章}
\subsection*{客观题答案}
1 B；2 错；3 C；4 A；5 B；6 对。
\begin{description}
\item[1--5（评分要点）] 控制器负责协议、服务负责用例事务、Repository负责持久化、数据库负责约束；同类内部自调用绕过代理，可拆分 Bean 或通过代理入口调用；接口写方法、路径、请求校验和 2xx/4xx/5xx；事务事件适合进程内提交后动作，Kafka 适合跨服务可靠传播；JWT 密钥轮换保留重叠窗口、刷新令牌轮换并可撤销。
\item[实践1] 验收必须包含正常、参数错误、无权限、资源不存在、重复事件和审计日志，测试可重复运行。
\end{description}

\section{第6章}
\subsection*{客观题答案}
1 B；2 C；3 错；4 对；5 对；6 错。
\begin{description}
\item[1--6（评分要点）] 说明 $M,V,P$ 顺序和坐标空间；WMS 返回渲染图、WMTS 返回瓦片、WFS 返回要素；椭球高与正常高需记录基准和转换模型；特征/连接/像控/检查点分别用于匹配、连接、约束和独立验收；IFC 4.3 用原生实体或项目分类/属性集边界表达；五维映射需把物理、虚拟、服务、数据和连接与人工确认相连。
\item[实践1--3] GLB 题提交加载 Promise、坐标、包围盒和失败降级；检查点题给出 $m_x,m_y,m_h$ 计算和限差判断；验收表覆盖数据、场景、性能、权限、故障和复现证据。
\end{description}

\section{第7章}
\subsection*{客观题答案}
1 错；2 C；3 B；4 错；5 对。
\begin{description}
\item[1--5（评分要点）] 质量码决定连线、颜色和是否进入评分；坐标转换写清基准、单位和矩阵；图表按任务选型并解释误导风险；像素半径必须换算到世界单位；颜色之外提供形状、文字和图标。
\item[实践1--3] 1000 条曲线需 300 点窗口、缺测断线和原始数据入口；28 点映射提交坐标表和五点误差；联动题提交曲线事件、同一 assetId 定位、详情查询和网络/空数据失败测试。
\item[实践4--5] LTTB 题比较峰谷保留率并说明不可用于统计；无障碍题提交键盘焦点、触控、读屏、大屏窄屏和颜色对比的失败与修复证据。
\end{description}

\section{第8章}
\subsection*{客观题答案}
1 B；2 C；3 错；4 对；5 对。
\begin{description}
\item[1--6（评分要点）] 角色边界需含允许/禁止操作；事件链包含原始值、质量码、规则版本、预警、确认、工单和复盘；安全监测与“四预”按模型预演和预案闭环区分；页面契约写输入、输出、版本、权限和错误；闸门反算说明多解、约束排序和人工选择；AI 助手只调用白名单工具并保留确认和审批。
\item[实践1--5] Vue 题验收 28 点筛选、ECharts 缺测断线、三维定位和权限错误；质量接口覆盖格式/单位/范围/时序/重复及三值质量码；模型卡含适用范围和版本；Compose 题提交前后端、PostgreSQL、Redis、探针和配置；降级演练验证超时和场景失败后二维、曲线与工单仍可用。
\end{description}

\section{第9章}
\begin{description}
\item[1] 评分覆盖覆盖、功耗、时延、所有权、安全、运维六项；说明双链路或人工补录触发条件。
\item[2] 评分覆盖标准、开源实现、三项可复现基准、失败判据、退出条件和迁移成本。
\item[3] 评分覆盖数据分层、四类存储职责、查询路径、备份点和恢复验证。
\item[4] 评分覆盖输入快照、模型版本、质量码、人工复核、审批、审计及超时/越界降级。
\item[5] 评分覆盖迭代目标、可运行切片、安全与容量测试、故障演练、复盘、责任角色和证据格式。
\item[能力回顾与交付物] 以图/表串联第1--8章能力；提交业务场景、权限、接口、字典、架构记录和端到端切片；技术评估必须有退出条件，复盘计划必须有责任人和完成证据。
\end{description}
